% Standalone source; no external input, bibliography, figures, or generated files.
% Intentionally uncompiled, as requested. Established results and derived synthesis are distinguished.
\documentclass[11pt]{article}
\usepackage[T1]{fontenc}
\usepackage[utf8]{inputenc}
\usepackage[margin=1in]{geometry}
\usepackage{amsmath,amssymb,amsthm}
\usepackage[hidelinks]{hyperref}
\newtheorem{theorem}{Theorem}[section]
\newtheorem{proposition}[theorem]{Proposition}
\newtheorem{corollary}[theorem]{Corollary}
\newtheorem{lemma}[theorem]{Lemma}
\newtheorem{definition}[theorem]{Definition}
\theoremstyle{remark}
\newtheorem{remark}[theorem]{Remark}
\title{One operational construction for residual automata,
formal concepts, and topology}
\author{A standalone technical synthesis}
\date{4 November 2026}
\begin{document}
\maketitle
\section{What is being unified}
\label{nb:section}

The shared datum below does more than feed three unrelated constructions.
Its sequential operations extend uniquely to the concept lattice, and its
supported concept poset models the Dowker complex in a way that intertwines
the operations on their exact domains. The operations are generally
\emph{partial on simplices}, even when the underlying state transition is
total. This distinction supplies both a transport theorem and a sharp
obstruction to an unrestricted common action.

The core theorem is finite, deterministic and Boolean; a final section gives
the set-sized extension. A completed operational
presentation is required; an arbitrary finite sample does not supply the
missing continuation operations or identify an unknown global language.
The ingredients belong to established automata, FCA, Chu-space and Dowker
theory. The particular integration below is proved directly without a
priority claim.


The central result is not that three procedures accept the same table.
It is that compatible sequential operations have unique algebraic lifts,
and that the same closure comparison intertwines their maximal partial
geometric realizations, including exact domains and word composition
(Theorem~\ref{nb:partial-bridge}). The marked residual states sit inside
that algebraic completion. The supported part of the completion has the
same homotopy type as the native incidence complexes. These are precise
transport statements; the three unmarked outputs are not equivalent.

\paragraph{Status of the claims.}
Residual minimization, concept completion, concept adjunctions, and
FCA--Dowker comparisons are established mathematics. Proofs are included
because arrow variance, extrema, and action domains matter. The combined
partial-action comparison is presented as a direct synthesis of those
constructions, not as a claim of novelty. Reconstruction of an independently
given space is a separate conditional theorem in
Section~\ref{sec:reconstruction}, not a consequence of a binary table alone.
\section{Acted observations and predictive residual states}

For the finite bridge, let $H,E$ be finite sets and let $R\subseteq H\times E$. Write
\[
 r(h)=\{e:hRe\},\qquad C_e=\{h:hRe\}.
\]
For $U\subseteq H$ and $B\subseteq E$, put
\[
 U'=\{e:(\forall h\in U)\ hRe\},\qquad
 B'=\{h:(\forall e\in B)\ hRe\},\qquad c(U)=U''.
\]
The type of the argument determines the direction of the prime. Empty
intersections are included, so $\varnothing'=E$ on the object side and
$\varnothing'=H$ on the attribute side.

When discussing an automaton, require $H,E$ to be nonempty.
Fix a finite alphabet $\Sigma$. A finite acted incidence presentation supplies
$h_0\in H$, $e_0\in E$, and, for each $a\in\Sigma$, functions
\[
 u_a:H\longrightarrow H,\qquad v_a:E\longrightarrow E
\]
with the exact compatibility law
\begin{equation}
 u_a(h)Re\quad\Longleftrightarrow\quad hRv_a(e).
 \label{nb:adjoint-incidence}
\end{equation}
For $w=a_1\cdots a_k$, use
\[
 u_w=u_{a_k}\circ\cdots\circ u_{a_1},\qquad
 v_w=v_{a_1}\circ\cdots\circ v_{a_k}.
\]
Then $u_w(h)Re\Longleftrightarrow hRv_w(e)$,
$u_{wz}=u_z u_w$, and $v_{wz}=v_wv_z$.
Here $u_a$ represents appending an action to a history, while $v_a$
represents prepending that action to a continuation test. Finite sets of
literal words are not usually closed under these operations; $H$ and $E$
are finite handles or quotients with supplied operations.

Define $h\sim k$ when $r(h)=r(k)$ and let $Q=H/{\sim}$.
Equation~\eqref{nb:adjoint-incidence} makes
\[
 \delta_a([h])=[u_a(h)],\quad q_0=[h_0],\quad
 F=\{[h]:hRe_0\}
\]
a well-defined DFA. It recognizes
$L=\{w:u_w(h_0)Re_0\}$.

\begin{proposition}[Exactly when this is the residual DFA]
\label{nb:residual}
Assume every $h\in H$ is reachable from $h_0$, and that every test in $E$
has the form $v_w(e_0)$ for some word $w$. Then $Q$ is the minimal
reachable DFA for $L$. More generally, retaining reachability, the same conclusion holds whenever
the tests $\{v_w(e_0):w\in\Sigma^*\}$ separate every pair of distinct rows
of $R$.
\end{proposition}
\begin{proof}
For a state represented by $h$, its continuation language is
$L_h=\{w:hRv_w(e_0)\}$. Equal rows have equal continuation languages.
Under the separation hypothesis the converse holds. Reachability makes
these exactly the left residuals of $L$. Thus the displayed quotient is
the Myhill--Nerode quotient, not merely a quotient relative to a sampled
set of tests.
\end{proof}

Every finite DFA supplies such data: take its states for $H$, and for $E$
take the distinct subsets $u_w^{-1}(F)$ of states. Incidence is membership,
$e_0=F$, and $v_a(C)=u_a^{-1}(C)$. The set $E$ is finite and closed under
these inverse images. Restricting $H$ to reachable states and quotienting
its rows gives the minimal DFA. This constructs the finite presentation
from a known finite machine; it does not infer that machine from a finite
unstructured sample.


\subsection{The full language context and a finite operational presentation}
For a finite alphabet $\Sigma$ and any language $L\subseteq\Sigma^*$, the native
history--continuation context is
\[
 H=E=\Sigma^*,\qquad hRt\ \Longleftrightarrow\ ht\in L.
\]
Its row at $h$ is $h^{-1}L=\{t:ht\in L\}$, and the compatible operations
are $u_a(h)=ha$ and $v_a(t)=at$. The resulting residual machine has states
$Q_L=\{h^{-1}L:h\in\Sigma^*\}$, initial state $L$, transitions
$B\mapsto a^{-1}B$, and accepting states containing $\epsilon$.
Every state is reached by its defining history, and distinct states are
separated by a continuation in their symmetric difference. Consequently
this is the minimal reachable deterministic recognizer. Any other
reachable deterministic recognizer maps onto it by sending a state to its
future language; distinct residuals cannot be represented by the same state.
This proves minimality, not only well-definedness. The machine is finite
exactly when $L$ is regular.

\begin{proposition}[Finite reduction preserving the native structures]
For regular $L$, put
\[
 Q=\{h^{-1}L:h\in\Sigma^*\},\qquad
 E_0=\{C_t:t\in\Sigma^*\},\qquad C_t=\{q\in Q:t\in q\}.
\]
Let $\bar R\subseteq Q\times E_0$ be membership, $q\bar R C$ iff $q\in C$.
Then $|E_0|\leq2^{|Q|}$ and the operations
\[
 \bar u_a(q)=a^{-1}q,\qquad
 \bar v_a(C_t)=C_{at}=\bar u_a^{-1}[C_t]
\]
are well defined and form a finite acted presentation. The two quotient
maps $\alpha(h)=h^{-1}L$, $\beta(t)=C_t$ induce a full concept-lattice
isomorphism preserving support, Dowker homotopy equivalences on both
sides, and exact compatibility of quotient maps with all action domains.
\end{proposition}
\begin{proof}
The set $Q$ is finite and derivative-closed. The displayed inverse-image
formula makes $\bar v_a$ independent of the representative $t$. It also
gives $\bar u_a(q)\bar R C\Longleftrightarrow q\bar R\bar v_a(C)$.
Both quotient maps are surjective and
$hRt\Longleftrightarrow\alpha(h)\bar R\beta(t)$.
For subsets $U\subseteq Q$, $V\subseteq E_0$, surjectivity gives
\[
 (\alpha^{-1}[U])'_R=\beta^{-1}[U'_{\bar R}],\qquad
 (\beta^{-1}[V])'_R=\alpha^{-1}[V'_{\bar R}].
\]
Every original closed extent is a union of row classes and every closed
intent a union of column classes. Hence
$(U,V)\mapsto(\alpha^{-1}[U],\beta^{-1}[V])$ is a concept-lattice
isomorphism, inverse to taking direct images. Nonemptiness of both
coordinates is preserved by surjectivity.

For an original finite object set $S$, being a face is equivalent to
$\alpha[S]$ being a face, since test witnesses lift along $\beta$.
Choose one history representing each row. This gives a simplicial section
of $\alpha$; the composite section--quotient is contiguous to the identity,
because $S$ and all its chosen representatives have the same common tests.
Contiguous maps have homotopic realizations by affine interpolation inside
their common image simplex. Thus $\alpha$ is a homotopy equivalence;
the same argument applies to $\beta$ on the test nerve.
Finally $\alpha u_a=\bar u_a\alpha$ and
$\beta v_a=\bar v_a\beta$. For any word $w$, $u_w[S]$ is a face
exactly when $\bar u_w[\alpha[S]]$ is, proving equality of the inverse-image
action domains. The concept statement follows either from the displayed
prime identities or from closure and image compatibility. The initial
state is $L$, the acceptance test is $C_\epsilon$, and every test is
$\bar v_w(C_\epsilon)=C_w$. Representative sections need not commute
with the operations; the quotient maps do.
\end{proof}
Thus the finite bridge applies to every regular language without assuming
that a closed set of literal histories is finite.

\subsection{What finite observation and learning do not supply}
The finite subsets of words used in an observation table need not be
closed under $u_a$ or $v_a$. Row equality relative to sampled tests is not
necessarily a predictive congruence. For $L=\{a\}$ over $\{a\}$ and the
single test $E_0=\{\epsilon\}$, histories $\epsilon$ and $aa$ have equal
empty rows, but successors $a$ and $aaa$ have different rows. Hence no
function on this row quotient, or on its object concepts, realizes the
append-$a$ operation on every history.

Active automata learning addresses this defect with structured access to
more information, rather than declaring the sample complete. Angluin's
observation-table construction uses closedness and consistency, with
membership and equivalence queries, to build and refine a conjectured
DFA~\cite{cg:angluin}. In the notation above, consistency demands that
row-identical histories have row-identical letter successors; closedness
ensures that the successor rows have representatives. A genuine positive
equivalence answer is global information. The construction here explains
what native structures become available once the required action and test
completeness are supplied; it is not a new learning algorithm or a proof
that an arbitrary sample determines its unknown source language.
\section{Concept completion and its canonical action}


The elementary polarity identity is
\[
 U\subseteq B'\quad\Longleftrightarrow\quad U\times B\subseteq R
 \quad\Longleftrightarrow\quad B\subseteq U'.
\]
Both prime maps reverse inclusion. Substituting $B=U'$ proves
$U\subseteq U''$; taking primes gives $U'''\subseteq U'$, and
extensivity on the test side gives the reverse inclusion. Thus
$U'''=U'$, and double prime is monotone, extensive and idempotent.
The same proof applies on the test side. These facts establish the closure
properties used below, including empty families.

Let $\mathcal L_R=\{A\subseteq H:c(A)=A\}$, ordered by inclusion.
This is the concept lattice in extent notation: the concept corresponding
to $A$ is $(A,A')$. Its joins are
\[
 \bigvee_i A_i=c\left(\bigcup_i A_i\right),
\]
and its meets are intersections. Put $\gamma(h)=c(\{h\})$.


\subsection{What the completion is free over, and what it is not}
Order intents by inclusion and write
\[
 \mathcal I_R=\{B\subseteq E:B''=B\}
 =\left\{\bigcap_{h\in U}r(h):U\subseteq H\right\}.
\]
Then $\mathcal L_R\cong\mathcal I_R^{\mathrm{op}}$; suppressing this
opposite order would reverse every preservation assertion below. Every
row is an intent by triple-prime equality. The map $h\mapsto\gamma(h)$
identifies equal rows and is injective on their quotient, and
\[
 A=\bigvee_{h\in A}\gamma(h)\quad(A\in\mathcal L_R).
\]
Thus residual states form a marked join-generating subset of the extent
lattice, equivalently a meet-generating subset of the intent lattice.
The inclusion is generally proper. If there are $n$ distinct rows, then
$n\leq |\mathcal L_R|\leq2^n$.

\begin{proposition}[A free presentation with incidence equations]
The powerset $\mathcal P(H)$, under union, is the free complete join
semilattice on $H$. The surjection $c:\mathcal P(H)\to\mathcal L_R$
is a complete-join homomorphism, where joins on the codomain are closed
unions. For a complete lattice $M$ and a map $b:H\to M$, there exists a
complete-join homomorphism $\bar b:\mathcal L_R\to M$ with
$\bar b(\gamma(h))=b(h)$ if and only if
\[
 c(U)=c(V)\ \Longrightarrow\
 \bigvee_{h\in U}b(h)=\bigvee_{h\in V}b(h)
 \quad(U,V\subseteq H).
\]
When it exists it is unique. Dually, a specified map $d$ from the distinct
rows into another intent lattice extends to a complete-meet map from the
source intent completion exactly when
\[
 \bigcap\mathcal U=\bigcap\mathcal V\ \Longrightarrow\
 \bigcap d[\mathcal U]=\bigcap d[\mathcal V]
\]
for all families of rows, including the empty family.
\end{proposition}
\begin{proof}
A complete-join map from $\mathcal P(H)$ must send $U$ to
$\bigvee_{h\in U}b(h)$, and this formula preserves arbitrary unions by
flattening the indexing families. The closure identity
$c(\bigcup_i c(U_i))=c(\bigcup_i U_i)$ proves that $c$ preserves the
stated joins. Its kernel consists exactly of the displayed closure
equations. The free map factors through its image precisely when it
respects those equations; the factor is unique by surjectivity, and its
join preservation follows from the same closed-union identity. The dual
argument defines the meet extension by
$\bigcap\mathcal U\mapsto\bigcap d[\mathcal U]$; equation preservation
makes it well-defined, and flattening intersections proves the claim.
\end{proof}
This is a free presentation \emph{modulo incidence equations}, not a free
complete lattice on the predictive-state set. Actual test subsets determine
which intersections coincide. The inclusion order of the row poset alone
does not retain all that information. In particular one cannot substitute
an unspecified ``free completion'' for this construction.

A Chu morphism $(f,g):R\longrightarrow S$, with
$S\subseteq K\times D$, consists of
\[
 f:H\longrightarrow K,\qquad g:D\longrightarrow E,
 \qquad f(h)Sd\ \Longleftrightarrow\ hRg(d).
\]
Let $c_S$ denote the closure for $S$.

\begin{theorem}[Canonical algebraic lift]
\label{nb:lattice-lift}
A Chu morphism induces an adjoint pair
\[
 f_!: \mathcal L_R\rightleftarrows\mathcal L_S:f^*,
 \qquad f_!(A)=c_S(f[A]),\qquad f^*(B)=f^{-1}[B],
 \qquad f_!\dashv f^*.
\]
In particular, $f_!$ preserves all joins, including the empty join.
It is the unique join-preserving map satisfying
\begin{equation}
 f_!\gamma(h)=\gamma_S(f(h)).
 \label{nb:generator-equivariance}
\end{equation}
These lifts preserve identity and composition. Their intents obey
\begin{equation}
 (f_!A)'=g^{-1}[A'].
 \label{nb:intent-transport}
\end{equation}
Consequently every $u_a$ has a canonical join-preserving extension
$D_a=(u_a)_!$ to the concept lattice, with
$D_w=D_{a_k}\cdots D_{a_1}$ and $D_{wz}=D_zD_w$.
The map $[h]\mapsto\gamma(h)$ embeds $Q$ as a distinguished,
join-generating subset, stable under all $D_a$.
\end{theorem}
\begin{proof}
The inverse image under $f$ of the target extent $B_0'$ is
$(g[B_0])'$, hence is closed. Thus inverse image restricts to closed
extents. For closed $A,B$,
\[
 c_S(f[A])\subseteq B\ \Longleftrightarrow\ f[A]\subseteq B
 \ \Longleftrightarrow\ A\subseteq f^{-1}[B],
\]
which proves the adjunction. The same inverse-image property gives
$f[c(U)]\subseteq c_S(f[U])$ for every $U\subseteq H$, hence
\begin{equation}
 c_S(f[c(U)])=c_S(f[U]).
 \label{nb:closure-naturality}
\end{equation}
This proves~\eqref{nb:generator-equivariance}. Every closed extent is
$\bigvee_{h\in A}\gamma(h)$, so that equation forces uniqueness.
Composition follows from~\eqref{nb:closure-naturality}, or from composition
of the right adjoints. Finally
$(f[A])'=g^{-1}[A']$, and closing a set does not change its intent,
which proves~\eqref{nb:intent-transport}.
Since $\gamma(h)'=r(h)$, two object concepts coincide exactly when the
two rows coincide. This gives the claimed embedding of $Q$.
\end{proof}


\paragraph{Exactly what happens to implications.}
An attribute implication $B\Rightarrow e$ holds in a context precisely
when $e\in B''$, equivalently $B'\subseteq C_e$. For a Chu morphism
$(f,g):R\to S$, the map on closed intents is $I\mapsto g^{-1}[I]$.
It preserves inclusions and arbitrary intersections, so it preserves
semantic inequalities between intersections of closed intents. This is
not a claim about arbitrary substitutions in a syntactic implication
basis. There is also a precise syntactic rule in the opposite direction:
if $B\Rightarrow d$ holds in $S$, then
$g[B]\Rightarrow g(d)$ holds in $R$. Indeed a source object satisfying
all $g(b)$ has image under $f$ satisfying all $b$, hence $d$, and so
it satisfies $g(d)$. For source-to-target transport, inverse images of
an arbitrary premise can discard attributes outside $\operatorname{im}(g)$;
no converse rule without further hypotheses is asserted.

\begin{remark}[A completion, not additional native histories]
The original automaton is recovered on the marked generators
$\gamma(Q)$, not by declaring every lattice element to be an observed
history. A concept may be a new join of several such generators.
Acceptance is recovered by
$\gamma(h)\subseteq C_{e_0}$, and the initial point is $\gamma(h_0)$.
More generally the full row-reduced incidence is recovered from the
marked lattice using
$hRe\Longleftrightarrow\gamma(h)\subseteq C_e$.
Forgetting the object and test markings is a further abstraction, not an
equivalence of all the native structures.
\end{remark}

\begin{proposition}[Exact extension obstruction without a Chu witness]
\label{nb:extension-test}
For an arbitrary function $f:H\to K$, the following are equivalent:
\begin{enumerate}
\item There is a join-preserving $F:\mathcal L_R\to\mathcal L_S$ with
$F\gamma(h)=\gamma_S(f(h))$ for every $h$.
\item $f[c(U)]\subseteq c_S(f[U])$ for all $U\subseteq H$.
\item For each target test $d\in D$, the set $f^{-1}[C_d]$ is a closed
extent of $R$.
\item Whenever $c(U)=c(V)$, one has $c_S(f[U])=c_S(f[V])$.
\end{enumerate}
When these hold, $F(A)=c_S(f[A])$ and is unique. A Chu witness $g$
is stronger than these conditions: it requires each inverse test extent
to be an \emph{existing single test column}, not merely an intersection
of columns.
\end{proposition}
\begin{proof}
A join extension must send
$c(U)=\bigvee_{h\in U}\gamma(h)$ to $c_S(f[U])$, giving (4) and
uniqueness. Taking $V=c(U)$ gives (2). Condition (2) gives
closure-naturality~\eqref{nb:closure-naturality}, from which the stated
formula preserves joins and generators, giving (1).
For completeness, (3) implies that inverse images of all closed extents
are closed, since each is an intersection of test columns. Then
$f^{-1}[c_S(f[U])]$ is closed and contains $U$, so it contains $c(U)$,
which gives (2). Conversely, let $B$ be a closed target extent and put
$U=f^{-1}[B]$. Condition (2) gives
$f[c(U)]\subseteq c_S(f[U])\subseteq B$. Thus $c(U)\subseteq U$,
and extensivity gives equality. This proves (3).
\end{proof}

\subsection{The language specialization and the universal automaton}
We now write $\widehat D_a$ for the action on intents, to distinguish it
from the extent action $D_a$. The pair notation $F_a(A,A')$ expresses the
same concept action as $D_a(A)$; extent and intent orders are opposite.
\paragraph{Theorem (canonical deterministic extension).}
In the full language context, every letter acts on the entire intent lattice by
\[
 \widehat D_a(B)=a^{-1}B.
 
\]
This is an endofunction preserving arbitrary intersections, including the
empty intersection. It is the unique complete-meet-preserving extension of
the derivative action on the distinct residual rows $Q_L$. Dually, on the standard concept lattice it is
\[
 F_a(A,B)=\bigl((Aa)'',a^{-1}B\bigr),\qquad
 Aa=\{ha:h\in A\},
 
\]
and preserves arbitrary joins. It satisfies
$F_a\gamma(h)=\gamma(ha)$. More generally $\widehat D_w(B)=w^{-1}B$ and
$\widehat D_v\widehat D_u=\widehat D_{uv}$; equivalently $F_vF_u=F_{uv}$.
From the initial concept $\gamma(\epsilon)$, the reachable concepts are
exactly the image of $Q_L$.
\begin{proof}
For every $A\subseteq\Sigma^*$,
\[
 a^{-1}\left(\bigcap_{h\in A}h^{-1}L\right)
 =\bigcap_{h\in A}(ha)^{-1}L=(Aa)'.
 
\]
Thus the derivative is an intent and commutes with all intersections.
All intents are intersections of residuals, so an extension with this
preservation property is forced by the displayed intersection identity.
The extent belonging to $(Aa)'$ is $(Aa)''$, proving the displayed
concept-pair formula. The remaining identities follow directly from
concatenation and the initial state.
\end{proof}
A stronger blanket assertion that ``the concept lattice has no deterministic
transition'' is therefore false for the native full language context.
It has a canonical transition; what can fail is preservation of the other
lattice operation, restriction to a selected subposet, or descent after
truncating the tests.

These constructions already occur in automata theory. A concept $(A,B)$ of
the full language context is exactly a maximal factorization $AB\subseteq L$. These are
the states of the universal automaton. The identification with intersections
of left quotients and the universal-automaton definition are explicit in
Lombardy--Sakarovitch, Proposition 2.1 and Definition 2.4
\cite{cg:universal}. For regular $L$, equipping these intents with the
deterministic derivative gives the canonical finite meet automaton of
Kl\'{\i}ma--Pol\'{a}k, Section 4 \cite{cg:meet}. For arbitrary $L$, the
theorem gives its complete-meet analogue using arbitrary intersections.

The nondeterministic universal automaton and this deterministic
complete-meet recognizer have the same underlying concept states but
different transition structures.
For concepts $c=(A,B)$ and $d=(C,D)$ its edge relation is
\[
 c\xrightarrow{a}_{\mathrm{UA}}d
 \Longleftrightarrow AaD\subseteq L
 \Longleftrightarrow D\subseteq a^{-1}B
 \Longleftrightarrow F_a(c)\leq d.
 
\]
Its initial states form the upset of $\gamma(\epsilon)$; its final states
satisfy $\epsilon\in B$. Thus $F_a(c)$ is the least universal-automaton
$a$-successor in standard concept order. No claim identifying the two
\emph{automata}, or their reachability structures, follows from identifying
their state sets.

\subsection{A forced failure of the other lattice operation}
Take $\Sigma=\{a,b\}$ and $L=\{\epsilon,a,ba\}$, with all tests.
The distinct residuals and all intents are respectively
\[
 Q_L=\{\varnothing,\{\epsilon\},\{a\},L\},\qquad
 \mathcal I_L=Q_L\cup\{\Sigma^*\}.
\]
Indeed the residuals after $\epsilon,a,b$ are $L,\{\epsilon\},\{a\}$;
every other residual is either $\{\epsilon\}$ or $\varnothing$. Their
intersections produce no further sets except the empty intersection
$\Sigma^*$. In the intent lattice,
\[
 \{\epsilon\}\vee\{a\}=L,\qquad
 \widehat D_b(L)=\{a\},\qquad \widehat D_b(\{\epsilon\})=\widehat D_b(\{a\})=\varnothing.
\]
All three inputs in this join equation are residual states. Consequently
\emph{no binary-join-preserving endofunction} of the intent lattice can
extend the residual $b$-transition. In particular no lattice homomorphism,
even one not required to preserve extrema, can extend it. The canonical
complete-meet-preserving derivative $\widehat D_b$ nevertheless exists, as proved
above. Dually, the standard concept action preserves joins but cannot be
made meet-preserving while agreeing with the residual transitions.
This example also separates the FCA completion from the
Dedekind--MacNeille completion of the residual poset: $Q_L$ is already a
four-element complete lattice, whereas $\mathcal I_L$ has a fifth element.

\section{The supported core and its Dowker topology}

Let $K_R$ be the object Dowker complex: its nonempty simplices are the
nonempty $U\subseteq H$ with $U'\ne\varnothing$. Objects with empty rows
are not vertices. The empty simplex is included by convention. Write
$\mathcal F(K_R)$ for its poset of nonempty faces. Define
\[
 \mathcal C_R=
 \{A\in\mathcal L_R:A\ne\varnothing,\ A'\ne\varnothing\}.
\]
This is the supported concept core. Let $\Delta P$ be the order complex
of a finite poset $P$, and write $|P|:=|\Delta P|$.

\begin{lemma}[Order homotopy and closure retracts]\label{cg:order-homotopy}
If monotone maps $u,v:P\to Q$ satisfy $u(p)\leq v(p)$ for every $p$, then
$|u|$ and $|v|$ are homotopic, relative to the subposet on which they agree.
If $c:P\to P$ is monotone, extensive, and idempotent, then
$|c(P)|\subseteq|P|$ is a strong deformation retract.
\end{lemma}
\begin{proof}
Define $H:P\times\{0<1\}\to Q$ by $H(p,0)=u(p)$ and $H(p,1)=v(p)$.
It is monotone: when $p\leq p'$, the only mixed comparison needed is
$u(p)\leq u(p')\leq v(p')$.
The order complex of the product supplies the usual staircase triangulation
of $|P|\times[0,1]$. Realizing $H$ therefore gives the required homotopy.
It is stationary on a subposet where $u=v$. For the second assertion take
$u=\mathrm{id}_P$ and $v=c$. If $i:c(P)\hookrightarrow P$ is inclusion,
then $c i=\mathrm{id}$, and the homotopy from $\mathrm{id}_P$ to $i c$
fixes $c(P)$ pointwise. The statement also holds when $P$ is empty. In this note all such posets are finite unless explicitly stated otherwise.
\end{proof}

\begin{proposition}[Concept closure realizes the Dowker comparison]
\label{nb:dowker-core}
The maps
\[
 r_R:\mathcal F(K_R)\longrightarrow\mathcal C_R,
 \quad U\longmapsto c(U),
 \qquad j_R:\mathcal C_R\hookrightarrow\mathcal F(K_R)
\]
satisfy $r_Rj_R=1$ and $1\leq j_Rr_R$. They induce homotopy
equivalences
\[
 |K_R|\simeq|\Delta\mathcal C_R|\simeq|K_{R^{\mathsf T}}|.
\]
Here the second comparison uses the order-reversing bijection
$(A,A')\leftrightarrow(A',A)$ and the corresponding closure on tests.
\end{proposition}
\begin{proof}
A nonempty simplex $U$ has nonempty closure and
$(c(U))'=U'\ne\varnothing$, so closure stays within the face poset.
It is an extensive idempotent monotone map whose fixed points are exactly
$\mathcal C_R$. The displayed poset inequalities yield homotopies of
order-complex maps and a retraction. The order complex of the face poset
is the barycentric subdivision of $K_R$. Apply the same argument on the
test side, where the supported concepts have the opposite order.
\end{proof}

Deleting the bottom and top elements indiscriminately is incorrect.
A bottom or top concept must be retained exactly when both its coordinates
are nonempty. If a supported bottom or top is present, the supported
poset is contractible. The support definition also handles empty carriers.


The full lattice has a least element and hence a contractible order
complex: its vertex for the least element is a cone point. More exactly,
writing $\bot=(E',E)$ and $\top=(H,H')$, the exact formula is
\[
 \mathcal C_R=\mathcal L_R\setminus
 \bigl(\{\bot:E'=\varnothing\}\cup\{\top:H'=\varnothing\}\bigr).
\]
The braces mean the indicated singleton when its condition holds, and the
empty set otherwise. Indeed an empty-extent concept must be bottom and
an empty-intent concept must be top. Both removal conditions apply even
when bottom and top coincide. The abstract full lattice together with
these two emptiness flags determines the supported poset and thus the
homotopy type; without these flags it need not do so.
For the $1\times1$ all-one relation the one concept has both coordinates
nonempty. Its nerve is a point; deleting its top and bottom would instead
give the empty space.

\section{Maximal-domain partial actions and the comparison theorem}

For a Chu morphism $(f,g):R\to S$, define a subcomplex and a subposet by
\begin{align}
 K_f&=\{\varnothing\}\cup
 \{\varnothing\ne U\subseteq H:U'\cap\operatorname{im}(g)\ne\varnothing\},
 \label{nb:partial-complex}\\
 \mathcal C_f&=
 \{A\in\mathcal L_R:A\ne\varnothing,
                    \ A'\cap\operatorname{im}(g)\ne\varnothing\}.
 \label{nb:partial-core}
\end{align}
Equivalently $K_f$ is the Dowker complex of the relation
$hR_f d\Longleftrightarrow f(h)Sd$, or of $R$ restricted to the tests
$\operatorname{im}(g)$. It is a subcomplex of $K_R$.

\begin{theorem}[One coherent partial action in the two native languages]
\label{nb:partial-bridge}
Every Chu morphism has the following properties.
\begin{enumerate}
\item $K_f$ is the \emph{largest} subcomplex of $K_R$ on which the vertex
function $f$ is a simplicial map into $K_S$. The restriction
$f_!:\mathcal C_f\to\mathcal C_S$ is the largest restriction of $f_!$
to a subposet of $\mathcal C_R$ whose values lie in $\mathcal C_S$.
\item Original source closure restricts to a retraction
$r_f:\mathcal F(K_f)\to\mathcal C_f$ with
$r_fj_f=1$ and $1\leq j_fr_f$. In particular it is a homotopy
equivalence after taking order complexes. This uses closure in $R$,
not a silently substituted closure in $R_f$.
\item The closure retractions commute strictly with both arrows in the
spans
\[
 \mathcal F(K_R)\ \longleftarrow\ \mathcal F(K_f)
       \ \xrightarrow{\ U\mapsto f[U]\ }\ \mathcal F(K_S),
 \qquad
 \mathcal C_R\ \longleftarrow\ \mathcal C_f
       \ \xrightarrow{\ f_!\ }\ \mathcal C_S.
\]
Hence these spans are naturally, objectwise homotopy-equivalent.
\item If $(p,q):S\to T$ is a second Chu morphism, the domain of their
composite is exactly the composite partial-map domain:
\begin{align*}
 K_{pf}&=\{U\in K_f:f[U]\in K_p\},\\
 \mathcal C_{pf}&=\{A\in\mathcal C_f:f_!(A)\in\mathcal C_p\},\\
 (pf)_!&=p_!f_!.
\end{align*}
\end{enumerate}
Thus the bridge respects identities and all finite compositions, rather
than supplying unrelated homotopy equivalences for individual actions.
\end{theorem}
\begin{proof}
For every $U\subseteq H$,
\[
 (f[U])'=g^{-1}[U'].
\]
For nonempty $U$, its being nonempty is exactly
$U'\cap\operatorname{im}(g)\ne\varnothing$.
This proves maximality on nonempty faces; the empty face is included
separately by convention. Since $f[A]$ is nonempty for nonempty
$A$, equation~\eqref{nb:intent-transport} gives the corresponding
maximality on supported concepts.

The identity $(c(U))'=U'$ proves that original closure preserves and
reflects the domain. More precisely,
\[
 r_R^{-1}(\mathcal C_f)=\mathcal F(K_f).
\]
Its fixed points on that domain are precisely $\mathcal C_f$. The closure
retraction argument applies unchanged. Commutation with inclusion is
immediate, and commutation with the right arrow is exactly
$c_S(f[U])=f_!(c(U))$, proved in
\eqref{nb:closure-naturality}.

The composite Chu morphism is $(pf,gq)$. Its domain condition is
$U'\cap\operatorname{im}(gq)\ne\varnothing$. Equivalently
$g^{-1}[U']\cap\operatorname{im}(q)\ne\varnothing$, which says that
$f[U]$ lies in $K_p$ and already implies $U\in K_f$.
Replace $U'$ by $A'$ and use~\eqref{nb:intent-transport} to prove the
concept-domain formula. The lattice composition formula was proved in
Theorem~\ref{nb:lattice-lift}.
\end{proof}


\paragraph{Functorial content, with the domains retained.}
Let $\operatorname{ParSimp}$ have simplicial complexes as objects and arrows
$K\leftarrow D\to L$ with $D$ a specified subcomplex of $K$, the left
map its inclusion, and the right map simplicial. Compose by restricting
the first domain to the inverse image of the second domain. This is again
a subcomplex. Identity arrows have full domain. Theorem~\ref{nb:partial-bridge}
therefore gives a strict functor from finite Chu contexts to
$\operatorname{ParSimp}$, and an analogous functor to partial monotone maps
on supported concept posets with lower-set domains. Closure gives an
objectwise homotopy equivalence of the corresponding face-poset spans,
on each full object and each domain, with strictly commuting squares.
The inverse inclusions of closed extents need not commute with the actions:
$f[A]$ need not already be closed. Thus no strictly equivariant inverse or
unspecified coherent choice of homotopy inverses is claimed.

For an acted incidence presentation, let $K_w=K_{u_w}$ and
$\mathcal C_w=\mathcal C_{u_w}$. The theorem gives, for every word,
\[
 K_R\longleftarrow K_w\xrightarrow{u_w}K_R,
 \qquad
 \mathcal C_R\longleftarrow\mathcal C_w\xrightarrow{D_w}\mathcal C_R,
\]
with
\[
 \mathcal C_w=\{A\ne\varnothing:A''=A,
                  \ A'\cap\operatorname{im}(v_w)\ne\varnothing\}.
\]
The exact word laws are
\[
 K_{wz}=\{U\in K_w:u_w[U]\in K_z\},\qquad
 \mathcal C_{wz}=\{A\in\mathcal C_w:D_w(A)\in\mathcal C_z\}.
\]
In particular $K_{wz}\subseteq K_w$: a simplex that has lost every
common test cannot regain one by a later action under this adjunction law.
This is a partial action with domains that are \emph{subcomplexes}, not
merely subsets of vertices. All its vertices may remain valid while one
of its edges ceases to be valid.

More precisely, let $\operatorname{Par}(P)$ consist of partially defined
monotone self-maps of a poset $P$ whose domains are lower subsets, with
ordinary partial-function composition. Each domain
$\mathcal F(K_w)$ is lower in $\mathcal F(K_R)$, and each
$\mathcal C_w$ is lower in $\mathcal C_R$. The formulas above give
\emph{strict} right monoid actions of $\Sigma^*$ in these two partial-map
monoids: the empty word is the everywhere-defined identity, and the
partial map for $wz$ is exactly the partial map for $w$ followed by that
for $z$, including equality of domains. Thus ``partial action'' here is
not an informal statement or merely a one-sided inclusion of domains.

Strict naturality above is asserted at the face-poset level. When a
simplicial map identifies vertices, the usual geometric identification
of a barycentric subdivision with its original complex is natural only
up to homotopy in general. The corresponding comparison of the geometric
spans therefore lives naturally in the homotopy category. No global
endomorphism of $H_*(K_R)$ is asserted when the action has a proper domain.

\begin{corollary}[Exact criterion for total topological transport]
\label{nb:total-criterion}
For a Chu morphism $(f,g):R\to S$, the following are equivalent:
\begin{enumerate}
\item $f$ induces a simplicial map $K_R\to K_S$ on all faces.
\item $f_!$ sends every supported concept into a supported concept.
\item Every nonempty source column $C_e$ is contained in some
$C_{g(d)}$, with $d\in D$.
\item $K_f=K_R$, or equivalently $\mathcal C_f=\mathcal C_R$.
\end{enumerate}
Surjectivity of $g$ is sufficient but not necessary. For a sequential
presentation, all words act totally on the nerve if and only if this
criterion holds for every letter.
\end{corollary}
\begin{proof}
Every nonempty simplex is contained in a source column. A nonempty column itself
is a simplex and is closed. It belongs to $K_f$ precisely when it is
contained in $C_{g(d)}$ for some $d$, giving the equivalences. Total
letter actions compose to total word actions; the converse includes
one-letter words.
\end{proof}

The minimal test in this criterion is column coverage, not literal
surjectivity. Repeated or dominated tests can make the latter unnecessary.
When $D$ is nonempty, one can choose a forward witness $\lambda(e)\in D$ with
$C_e\subseteq C_{g(\lambda(e))}$ whenever the criterion holds. This
turns the Chu morphism into a covariant incidence-preserving pair
$(f,\lambda)$ for purposes of the nerve. The vertex map on the object
nerve is already determined by $f$ and does not depend on that choice.


\paragraph{No automatic linear action on homology.}
A word gives the pair of maps $H_k(K_w)\to H_k(K_R)$ induced by the
domain inclusion and by $u_w$. This is a span, not a canonical
endomorphism of $H_k(K_R)$. If the inclusion induces an isomorphism it
can be inverted, but that is an extra hypothesis. Homology also need not
preserve the pullbacks defining composition of partial maps, so these
spans cannot silently be promoted to a strict monoid representation in
linear spans. Maximal partial simplicial maps and the use of paired
homology maps in this situation are established tools in
Edelsbrunner--Jab\l o\'nski--Mrozek~\cite{cg:selfmap}.

\section{Sharp failures of stronger identifications}
\subsection{A live minimal automaton whose nerve action is not total}

Take a unary automaton with states $H=\{q,p,r\}$, initial state $q$,
accepting states $\{q,p\}$, and
\[
 u(q)=p,\qquad u(p)=r,\qquad u(r)=p.
\]
Its language is $\{\varepsilon\}\cup\{a^{2k+1}:k\geq0\}$.
Its three distinct suffix tests are represented by
$e_0=\varepsilon$, $e_1=a$, and $e_2=a^2$, with incidence
\[
 \begin{array}{c|ccc}
       &e_0&e_1&e_2\\\hline
 q&1&1&0\\
 p&1&0&1\\
 r&0&1&0
 \end{array}
 \qquad
 v(e_0)=e_1,\quad v(e_1)=e_2,\quad v(e_2)=e_1.
\]
All states are reachable, all have accepting continuations, and all rows
are different, so this is already the minimal DFA.

The nerve has edges $\{q,p\}$ and $\{q,r\}$. The action sends
$\{q,p\}$ to $\{p,r\}$, which is not a simplex. This is not a failure
caused by a dead vertex: both $p$ and $r$ remain vertices.
The closed extents are
\[
 \varnothing,\ \{q\},\ \{p\},\ \{q,p\},\ \{q,r\},\ H.
\]
The lattice action is perfectly well defined, but
$D(\{q,p\})=c(\{p,r\})=H$, whose intent is empty.
Its exact domain nerve has the edge $\{q,r\}$ together with the isolated
vertex $p$, and its exact concept domain is
$\{\{q\},\{p\},\{q,r\}\}$.
Thus a total native state transition, its total concept-lattice lift,
and its partial topological action are distinct, explicitly compatible
objects.

\subsection{The same lattice action can support a circle or a disk}

Let $H=E=\{1,2,3\}$ and use the complement of the identity matrix:
\[
 iRj\quad\Longleftrightarrow\quad i\ne j,
 \qquad
 R=\begin{pmatrix}0&1&1\\1&0&1\\1&1&0\end{pmatrix}.
\]
Every subset of $H$ is a closed extent, so $\mathcal L_R$ is the Boolean
lattice $B_3$. A nonempty subset has a common test exactly when it is a
proper subset of $H$. Therefore $K_R$ is the triangle boundary,
\[
 |K_R|\cong S^1,
 \qquad \mathcal C_R=B_3\setminus\{\varnothing,H\}.
\]
Now add one universally true test $t$. Intersecting with its column $H$
adds no closed extents, so the extent lattice remains exactly $B_3$.
But $H$ now has nonempty intent, and hence
\[
 K_{R^+}=\text{the full two-simplex},\qquad
 \mathcal C_{R^+}=B_3\setminus\{\varnothing\}.
\]
The first nerve has first Betti number $1$; the second is contractible
and has first Betti number $0$.

The example also preserves dynamics. Set $u=(123)$ on objects and
$v=(132)$ on the original tests, and set $v(t)=t$ after enrichment.
Then $u(i)Rj\Longleftrightarrow iRv(j)$, and both contexts induce the
same lattice automorphism $A\mapsto u[A]$. The nerve action rotates a
circle in the first context and a disk in the second. Thus even the
unmarked lattice \emph{together with its action} does not determine the
native topological output.

Taking state $1$ as initial and test $1$ as the acceptance test gives
the same minimal three-cycle DFA in both contexts. In the original
context all tests are suffix-generated. The new tautology test is not
suffix-generated from that acceptance test; it is an enrichment of the
test interface. It changes no state equivalence and no accepted word,
but changes common-test topology. This is precisely why an unannounced
test enrichment is not a harmless change of representation.


This example also distinguishes arrow types. The tautology-forgetting
map $R^+\to R$ with object function $\mathrm{id}_H$ and backward test
inclusion $E\hookrightarrow E\cup\{t\}$ is a Chu morphism. Its full
lattice lift is the identity on $B_3$, but its maximal simplicial domain
is precisely the boundary inside the filled triangle. Conversely the
enrichment $R\to R^+$ is a covariant incidence-preserving inclusion.
There is no Chu morphism in this direction with object function
$\mathrm{id}_H$: the new universal test would have to pull back to an
old test with column $H$, and none exists.

\subsection{One-bit instability of ordinary homology}
Even without adding a test, changing the single entry $R(1,1)$ from zero
to one creates a universally incident column and a universally incident
row. Both Dowker complexes change from the triangle boundary to a filled
triangle: pairwise attribute intersections were already nonempty, while
the triple intersection now becomes nonempty. Thus one bit changes
$\dim H_1$ from $1$ to $0$. This assertion concerns static topology; the
old cyclic action need not remain compatible after that change.
Both matrices do support identity Chu actions, so requiring an action
alone does not eliminate the instability.

More generally, the complement of the $n\times n$ identity, for $n\geq3$,
has object nerve the boundary of an $(n-1)$-simplex. Flipping a single
diagonal entry fills that simplex and changes its $(n-2)$nd Betti number
from $1$ to $0$, at normalized entrywise Hamming distance $1/n^2$.
This exact incidence theorem therefore gives no uniform perturbation
stability for ordinary homology under that metric. A noisy application
needs its own noise assumptions and, if a filtration is introduced, a
separately justified persistence or other stability statement.

\subsection{Different native morphism classes really differ}

A covariant relation morphism $(f,\lambda)$ needs only
$hRe\Rightarrow f(h)S\lambda(e)$. It always induces an object nerve map.
It need not induce a generator-compatible join map of FCA lattices.
For example, use the $3\times3$ identity incidence and put
$f(1)=f(2)=1$, $f(3)=2$, $\lambda=f$. The pair preserves incidence.
Its concept lattice has three atoms and their common top. If a join map
extended $f$ on the atoms, the join of atoms $1,2$ would map to atom $1$,
while the identical top, expressed as the join of atoms $1,3$, would map
to top. Such an extension is impossible.

Conversely, a Chu morphism always has its FCA adjoint but need not have
a total nerve map, as the unary example shows. Thus neither native arrow
class can be silently substituted for the other. The exact-domain
construction of Theorem~\ref{nb:partial-bridge} handles all Chu arrows;
Corollary~\ref{nb:total-criterion} identifies the subcategory on which the
topological action is total.

The finite-reduction proposition also shows precisely why duplicate rows
and columns are benign: their quotient preserves support and the lattice,
and is a nerve homotopy equivalence commuting with the actions. It is
not an isomorphism of the original presented simplicial complexes.

\section{Covariant topology and the distinct Chu variance}

Let $\mathsf{Inc}_{\mathrm{fin}}$ have finite relations as objects and
morphisms
\[
 (f,g):(X,Y,R)\longrightarrow(\bar X,\bar Y,S),\qquad
 xRy\ \Longrightarrow\ f(x)Sg(y),
\]
where $f:X\to\bar X$ and $g:Y\to\bar Y$. Composition is componentwise.
This is a category of incidence structures, not the category whose objects
are sets and whose arrows are relations.


For use in the naturality proof, define on finite nonempty face posets
$d_R(\sigma)=\sigma'_R$ from $F(D_X(R))$ to
$F(D_Y(R))^{\mathrm{op}}$. It is well-defined because $\sigma'$ is
nonempty and has witness set $\sigma''\supseteq\sigma$. The symmetric
prime map has composites equal to closure. The order-homotopy lemma
therefore makes their realizations inverse homotopy equivalences. We use
$D_X,D_Y$ for the two Dowker complexes in this section, and $F$ for the
nonempty-face-poset construction.
\begin{proposition}[Strict Dowker functors and natural homotopy duality]
\label{cg:covariant-natural}
The assignments $D_X$ and $D_Y$ are strict functors from
$\mathsf{Inc}_{\mathrm{fin}}$ to finite simplicial complexes, with maps
induced by $f$ and $g$. Their geometric realizations are naturally
isomorphic after passage to the homotopy category of spaces.
Consequently their homology functors are naturally isomorphic, for every
fixed coefficient ring and every nonnegative degree.
\end{proposition}
\begin{proof}
If $y$ witnesses $\sigma$, then $g(y)$ witnesses $f(\sigma)$, so $f$
induces a simplicial map; the other side is symmetric. Images of subsets
respect identities and composition exactly.
On face posets the inequality
\[
 g(\sigma'_R)\subseteq(f(\sigma))'_S
\]
holds because every pair in $\sigma\times\sigma'_R$ is an incidence.
Viewed in $F(D_Y(S))^{\mathrm{op}}$, it says
\[
 d_S\,F(D_X(f,g))\ \leq\ F(D_Y(f,g))^{\mathrm{op}}\,d_R.
\]
Lemma~\ref{cg:order-homotopy} gives the homotopy commutativity of the
naturality square. Each $d_R$ is a homotopy equivalence by the preceding prime-map argument. Its homotopy class therefore
is a natural isomorphism. For a general simplicial map that identifies vertices, the canonical
barycentric-subdivision homeomorphism need not commute with that map
strictly. It commutes up to homotopy: on each subdivided simplex lying in
an original face $\sigma$, both composites land in the simplex $f(\sigma)$,
so straight-line interpolation in that simplex gives a compatible homotopy.
This is precisely the compatibility required in the homotopy category.
Homotopy-invariant homology sends the homotopies to equal maps and the
equivalences to isomorphisms.
\end{proof}

\subsection{A strict natural span using rectangles}

The preceding homotopy-natural equivalence already suffices for persistence.
For a strict natural span before taking homotopy classes, the rectangle
complex of Brun--Salbu~\cite{nb:brun-salbu} is useful:
\[
 E(R)=\{\varnothing\}\cup
 \{\varnothing\neq Q\subseteq R:
           \pi_X(Q)\times\pi_Y(Q)\subseteq R\}.
\]
Here the vertices are incidences, and $\pi_X,\pi_Y$ are the coordinate
projections. The map $(f,g)$ sends $(x,y)$ to $(f(x),g(y))$, defining a
strict functor $E$. Projections give a strictly natural span
\[
 D_X(R)\ \xleftarrow{\pi_X}\ E(R)\ \xrightarrow{\pi_Y}\ D_Y(R).
\]
Both projection maps are homotopy equivalences. We give a proof to avoid
mistaking a strict span for a strict chosen inverse.

\begin{lemma}[Finite-poset fiber lemma]\label{cg:finite-fiber}
Let $P,Q$ be finite posets. If $p:P\to Q$ is monotone and every
$P_q=\{a\in P:p(a)\leq q\}$ has contractible realization, then $|p|$
is a homotopy equivalence.
\end{lemma}
\begin{proof}
Form the poset $M=P\sqcup Q$ retaining both orders and adding
$a\leq q$ exactly when $p(a)\leq q$. The map $r:M\to Q$ equal to $p$
on $P$ and the identity on $Q$ is monotone, and
$\mathrm{id}_M\leq i_Qr$. Thus $i_Q:Q\hookrightarrow M$ is a homotopy
equivalence. Remove elements of $Q$ from $M$ in a linear extension of $Q$
from smaller to larger. At the removal of $q$, its strict lower poset is
exactly $P_q$, since the earlier elements of $Q$ have already been removed.
The link of the vertex $q$ in the order complex is the join of this lower
poset complex with the strict upper poset complex; it is contractible
because $|P_q|$ is nonempty and contractible. Deleting a vertex with
contractible link is a homotopy equivalence: the old complex is the union
of the deleted complex with the cone on the link, attached along that link;
a contractible link and its cone are homotopy equivalent, and the
subcomplex inclusion is a cofibration, so gluing preserves this homotopy
equivalence. Induction shows $i_P:P\hookrightarrow M$ is a homotopy
equivalence. Since $p=r i_P$, so is $p$.
If $Q$ is empty then $P$ is empty and the result is immediate.
\end{proof}

\begin{proposition}[Rectangle projections]\label{cg:rectangle}
Both $|\pi_X|$ and $|\pi_Y|$ in the natural span are homotopy equivalences.
\end{proposition}
\begin{proof}
Consider the face-poset map $p_X:F(E(R))\to F(D_X(R))$ sending $Q$ to
$\pi_X(Q)$. Fix a nonempty face $\sigma$ of $D_X(R)$ and a witness
$y\in\sigma'$. Its lower fiber consists of nonempty faces $Q$ with
$\pi_X(Q)\subseteq\sigma$. Put $A=\pi_X(Q)$ and $B=\pi_Y(Q)$.
Within this fiber the following are monotone maps and pointwise comparisons:
\[
 Q\ \subseteq\ A\times(B\cup\{y\})
 \ \supseteq\ A\times\{y\}
 \ \subseteq\ \sigma\times\{y\}.
\]
Every displayed set is a nonempty face of $E(R)$, because $A\times B$
lies in $R$ and $y$ witnesses every element of $\sigma$. The last map is
constant. Repeated application of the order-homotopy lemma contracts the
fiber. Apply Lemma~\ref{cg:finite-fiber} and barycentric subdivision.
The other projection is symmetric. If $R$ is empty, all spaces are empty.
\end{proof}

Here $E(R)$ denotes a complex, distinct from the test set $E$ used earlier.
The facets of $E(R)$ are exactly the nonempty maximal rectangles, hence
exactly the supported concepts viewed as subsets $A\times B$ of $R$.
This is another literal use of the same data. However, a choice of homotopy
inverse to $\pi_X$ is generally not strictly natural. It is the span that
is strictly natural; the induced isomorphism is natural in the homotopy
category.


\paragraph{Ordinary incidence maps do not repair the algebraic variance.}
The counterexample in the preceding section shows that a covariant
incidence pair need not admit the generator-compatible join extension.
The always-defined supported-concept assignment
$A\mapsto c_S(f[A])$ is then only lax under composition:
\[
 c_T(pf[A])\subseteq c_T\bigl(p[c_S(f[A])]\bigr).
\]
This follows from extensivity of $c_S$ and monotonicity. The inequality
becomes a homotopy after order realization, but need not be an equality
of lattice maps. By contrast, the Chu identity supplies the inverse-image
closure condition that forces equality. The covariant and contravariant
morphism classes serve different native purposes.

\paragraph{Compatible filtrations.}
For $R_s\subseteq R_t$ on fixed finite sets, both Dowker constructions
are filtrations and the natural homotopy equivalences induce isomorphic
persistence modules. This is the functorial strengthening associated with
Chowdhury--M\'emoli~\cite{cg:dowker-persistence}; the rectangle span is
another strict comparison before passing to homotopy. The concept posets
need not be nested subposets, and the naive concept maps need not compose
strictly; use their induced homology maps or the explicit spans instead.
Pointwise Betti numbers are insufficient: include a circle into a cylinder
in one two-stage filtration, but fill the original circle by a disk and
attach a new circle at a point in the other. Both stages in both cases have
the homotopy type of a circle, while the first $H_1$ structure map is an
isomorphism and the second is zero. A persistence comparison requires
compatible maps, not just pointwise equivalences.

\section{When an actual topological space is recovered}
\label{sec:reconstruction}
Let $Z$ be a space, let $\{U_e:e\in E\}$ be a finite family of regions,
and let $q:H\to Z$ be a finite observation map. Impose exact semantics
\[
 hRe\quad\Longleftrightarrow\quad q(h)\in U_e.
\]
For a nonempty set of labels $B$, write $U_B=\bigcap_{e\in B}U_e$.
The nerve $N(\mathcal U)$ has face $B$ precisely when $U_B\ne\varnothing$.
Direct substitution gives
\begin{equation}\label{cg:detected-intersections}
 B\in D_E(R)\quad\Longleftrightarrow\quad q(H)\cap U_B\ne\varnothing.
\end{equation}
Thus $D_E(R)\subseteq N(\mathcal U)$, with equality exactly when every
nonempty intersection is hit by a sample. Call this
\emph{intersection-complete sampling}. Since $E$ is finite, it suffices
to hit each maximal nerve face's intersection: every other nonempty
intersection contains one of those intersections. This condition is
weaker than observing every Boolean atom, but cannot be certified from
absent witnesses in the observed table alone.

\begin{theorem}[Conditional homotopy recovery]
Assume exact incidence semantics, $\bigcup_eU_e=Z$, and
intersection-complete sampling. Assume also either:
\begin{enumerate}
\item $Z=|A|$ for a finite simplicial complex $A$, each $U_e=|A_e|$
for a subcomplex $A_e$, and every intersection of the $A_e$ with
nonempty realization has contractible realization; or
\item $Z$ is paracompact Hausdorff, the $U_e$ are open, and every
nonempty finite intersection $U_B$ is contractible.
\end{enumerate}
Then
\[
 |K_R|\simeq|\Delta\mathcal C_R|\simeq|D_E(R)|
       =|N(\mathcal U)|\simeq Z.
\]
\end{theorem}
\begin{proof}
Equation~\eqref{cg:detected-intersections} proves equality of the observed
complex and the nerve. Proposition~\ref{nb:dowker-core} proves the first
two comparisons. We give the remaining argument fully in case (1).
Use the finite poset
\[
 P=\{(\sigma,J):\sigma\text{ a nonempty face of }A,\
 J\ne\varnothing,\ \sigma\in\bigcap_{j\in J}A_j\},
\]
ordered by inclusion of $\sigma$ and reverse inclusion of $J$.
Projection to $F(A)$ and projection to $F(N(\mathcal U))^{\mathrm{op}}$
are monotone. For a face $\sigma_0$ of $A$, choose an index $i$ with
$\sigma_0\in A_i$; such an index exists because the subcomplexes cover.
In the lower fiber over $\sigma_0$, the comparisons
\[
 (\sigma,J)\ \geq\ (\sigma,J\cup\{i\})\
 \leq\ (\sigma,\{i\})\ \leq\ (\sigma_0,\{i\})
\]
contract its realization. All pairs shown are valid because
$\sigma\subseteq\sigma_0\in A_i$. For a nerve face $J_0$, the lower
fiber of the other projection consists of pairs with $J\supseteq J_0$.
The map $(\sigma,J)\mapsto(\sigma,J_0)$ is an extensive retraction onto
$F(\bigcap_{j\in J_0}A_j)\times\{J_0\}$. Its realization is contractible
by assumption. The finite-poset fiber lemma
(Lemma~\ref{cg:finite-fiber}) makes both projections homotopy equivalences.
Barycentric subdivision and order reversal give
$|A|\simeq|N(\mathcal U)|$, as required.

In case (2), the last step is precisely the established good-open-cover
nerve theorem, rather than a new deduction from the matrix. A finite open
cover of a paracompact Hausdorff space is numerable; the hypotheses are
those of the good numerable open-cover case of the functorial nerve
theorem~\cite[Theorem 5.9]{cg:nerve}. Applying that theorem proves the
remaining equivalence. No assertion here depends on a nerve theorem for
arbitrary closed covers.
\end{proof}
The first proof makes the finite simplicial version independent of an
unmentioned general cover theorem. Useful additional established classes
include finite closed convex Euclidean covers; arbitrary finite good
closed covers are insufficient without tameness or cofibration hypotheses
\cite[Theorem 3.9 and Remark 5.12]{cg:nerve}.

\subsection{What the additional hypotheses exclude}
\begin{enumerate}
\item \emph{Unseen topology.} A single everywhere-true predicate gives
the same all-one matrix on a finite sample from an interval or a circle.
More generally, from any realization $(Z,U_e,q)$ form
$(Z\times W,U_e\times W,h\mapsto(q(h),w_0))$, for a pointed nonempty
space $W$. The incidence is unchanged, while the latent homotopy type
can change. Therefore no data-only reconstruction rule works for arbitrary
spaces and predicate interpretations.
\item \emph{Missing intersections.} A three-arc good open cover of a
circle can have all pairwise overlaps and no triple overlap. Observing
only an exclusive part of its first arc gives rows $(1,0,0)$ and an
observed point nerve, although the true nerve is a triangle boundary.
The failure is sampling completeness, even with valid cover topology.
\item \emph{A different topology generated by the same predicates.}
On $H=\{1,2,3\}$, use extents $\{1,2\},\{2,3\},\{1,3\}$. The topology
generated by these sets is discrete, because their intersections isolate
all three points. The attribute nerve is instead a circle. These are
both legitimate constructions from the data, but not the same topology.
\item \emph{Higher-order witnesses.} Nonempty pairwise overlaps do not
force a nonempty triple overlap. Replacing incidence by a pairwise graph
and taking its clique complex fills a triangle that the true nerve may
omit. Such replacement requires an additional flag condition.
\item \emph{Homotopy is not geometry.} A point and an interval have the
same homotopy type and the same one-predicate data. The theorem does not
recover homeomorphism type, metric, dimension, smooth structure, or
coordinates. If the regions cover only a subset of an ambient space,
only their union is the possible nerve-theorem target.
\end{enumerate}

\paragraph{Maps and persistence require semantic compatibility too.}
For another covered space, a continuous $F:Z\to Z'$ and a label map
$\lambda:E\to E'$ satisfying $F(U_e)\subseteq U'_{\lambda(e)}$ induce
a nerve map. If sample maps obey $Fq=q'f$, the observed covariant
incidence map agrees with that nerve map. Natural nerve comparison spans
then provide compatible topology comparisons. Applying this at every
parameter can identify persistence modules under the same sampling and
cover assumptions at every parameter. Separate pointwise equivalences do
not suffice, as the two-stage counterexample above demonstrates.
\section{Scope, limitations, and the contribution of the synthesis}

The common construction retains a finite incidence presentation, its
action adjunction, and the markings needed to recover native states and
tests. It produces the residual automaton under the stated completeness
hypotheses, a canonical generator-compatible lattice action, and a
coherent Dowker action by exact-domain spans. The closure maps intertwine
the last two constructions, including all action domains and word
compositions. This is a substantive common theorem.

It does not show that the three unmarked outputs are equivalent or that
one is recoverable from either of the others without retained structure.
It supplies neither an empirical discovery of the interface nor an
efficient uniform algorithm: concept lattices and nerves may have
exponentially many elements or faces. A finite sample of word labels
cannot certify global minimality without additional assumptions or
queries. If all sampled words have length at most $N$ and label zero,
both the empty language and $\{a^{N+1}\}$ fit the sample, while their
minimal complete DFAs have respectively $1$ and $N+3$ states.


\paragraph{Finite and infinite scopes.}
The polarity and complete-lattice action statements hold for arbitrary
sets. The closure deformation-retraction proof is finite: a finite face
may have infinite closure in an infinite context. Here is the precise
set-sized replacement. Assume that both carriers are sets, and define
Dowker complexes using finite nonempty faces. The monotone map
$r:\mathcal F(K_R)\to\mathcal C_R$, $S\mapsto S''$, is well defined
but need not have the earlier inclusion as an inverse. For a supported
concept $(A,B)$, the lower fiber $r^{-1}(\mathcal C_{R,\leq A})$ is
exactly the poset of finite nonempty subsets of $A$: since $A$ is closed,
$S''\subseteq A$ is equivalent to $S\subseteq A$. Every such subset is
a face, because the nonempty set $B=A'$ witnesses it. Choose $a_0\in A$.
The monotone comparisons
$S\subseteq S\cup\{a_0\}\supseteq\{a_0\}$ contract the fiber.
For a Chu morphism, restrict to $\mathcal F(K_f)\to\mathcal C_f$.
If $A\in\mathcal C_f$ and $S\subseteq A$, then $A'\subseteq S'$ and
$A'\cap\operatorname{im}(g)\ne\varnothing$, so $S$ lies in $K_f$.
The lower fiber is therefore the same contractible poset. All posets
involved are small because the carriers are sets. Quillen's Theorem A
in its poset form now makes these maps homotopy equivalences
\cite{cg:quillen}. The strict closure/action squares and exact domain
composition are algebraic identities and remain valid. As in the finite
case, passage from face posets to the original geometric complexes gives
homotopy naturality, not strict naturality under vertex identifications.
This is not the ill-typed inclusion of an infinite extent into a finite
face poset. Regular languages already admit the finite reduction above,
so no infinitary step is needed for that case.

\paragraph{Boolean predictions are not probability laws.}
For a binary Markov chain with $\Pr(1\mid0)=1/4$ and
$\Pr(1\mid1)=3/4$, every finite future word has positive probability
after either state. The Boolean possibility rows agree, although the
future distributions differ. Thus this Boolean theory does not identify
probabilistic predictive states. Adding rational probability-threshold
tests can recover equality of laws, but the update then uses
$\Pr(t\mid ha)=\Pr(at\mid h)/\Pr(a\mid h)$ when defined. The simple
column-shift identity and its algebraic consequences must be proved anew.
Likewise, a linear Hankel span uses a different free algebraic structure;
it is not obtained merely by relabeling the Boolean lattice construction.

The unification achieved here is a precise operational and topological
transport theorem, with an explicit failure boundary. It is neither a
full phenomenological theory of concept or scientific-law formation nor
an automatic account of how the history/test interface is discovered.
\section{Prior work and exact attribution}
\label{nb:prior}

The elementary ingredients are established and should be credited
separately. Myhill--Nerode theory supplies residual minimization.
Clark's discussion of the syntactic concept lattice explicitly compares
the prefix--suffix relation underlying deterministic quotients and
universal automata; the construction here is not a new discovery that
automata and concept lattices share incidence data
\cite{nb:clark}. Shen and Zhang prove concept-lattice functoriality for
infomorphisms in a substantially more general enriched setting
\cite{nb:shen-zhang}. Erne surveys several context-morphism categories and
their lattice adjunctions \cite{nb:erne}.

Gratus and Porter already restrict source attributes to the image of
the contravariant component of a Chu morphism to obtain a simplicial
map, and give surjectivity as a sufficient condition for an unrestricted
map. Their refinement discussion includes the column-containment
condition used above \cite{nb:gratus-porter}. Thus neither the restriction
to $\operatorname{im}(g)$ nor the coverage criterion should be presented
as new. Brun and Salbu explain prior FCA--Dowker comparisons and give a
functorial Dowker theorem for covariant relation morphisms
\cite{nb:brun-salbu}. Their arrow convention must not be confused with
the contravariant test component of a Chu morphism.

The claim made here is a checked synthesis: one closure retraction
intertwines all of the exact-domain action spans, with strict formulas
for composition; the examples identify precisely what totality and
unmarked-lattice recovery would falsely assert. An exhaustive novelty
assessment of this synthesis has not been performed.


The following distinctions further delimit the established components.
\begin{itemize}
\item Lombardy--Sakarovitch, Proposition 2.1 and Definition 2.4, identify
maximal language factorizations with intersections of residuals and define
the universal NFA~\cite{cg:universal}. Kl\'{\i}ma--Pol\'{a}k discuss the
corresponding deterministic meet automaton~\cite[Section 4]{cg:meet}.
The universal property of the universal NFA is existence of a morphism,
not generally uniqueness or terminality in an ordinary automaton category.
\item Clark's \emph{syntactic} concept lattice uses the different relation
$w\sim(u,v)$ iff $uwv\in L$. Its row equivalence is two-sided syntactic
congruence, unlike the prefix--suffix context's Myhill--Nerode relation.
It is relevant prior art, not a license to identify the two lattices.
\item For a small quantaloid and an enriched distributor, Shen--Zhang's
Proposition 4.1 gives an Isbell adjunction and Theorem 4.16 characterizes
its completion through sup/inf-dense maps~\cite{cg:quantaloid}. Boolean
discrete contexts recover FCA. Their Kan adjunction is a distinct
construction (Proposition 5.1); its relation to Isbell through negation
uses the Girard hypothesis (Lemma 5.12). These generalities do not, without
new hypotheses, identify a field-linear Hankel image with a Boolean nucleus.
\item Freund--Andreatta--Giavitto explicitly connect lattices and binary
relation topology, exclude universal rows/columns in their Definition 1,
and explain the need for extra marking if those are allowed
\cite{cg:freund}. Ayzenberg's Proposition 5.7 compares proper concept
parts with the two nerves under exclusions on full simplices and ghost
vertices~\cite{cg:ayzenberg}. The support convention above handles the
degenerate cases directly. The universal-test counterexample does not
contradict their qualified results.
\item Gratus--Porter's Proposition 1(ii) already supplies the restricted
Chu nerve map used here. Exact partial composition is a direct consequence
of its witness formula. The additional compatibility equations in this
note are proved explicitly; an exhaustive priority search for their exact
combined formulation has not been performed.
\item Green--Joslyn--Myers--Rawson--Robinson's 2025 preprint relates
hypergraph intersection structures, concept lattices, and Dowker cosheaves
\cite{cg:green}. It is further close prior art to any broad novelty claim,
not a foundation on which the proofs here rely.
\end{itemize}
The proofs, counterexamples, and exact preservation requirements are the
substance of this note. No claim of new mathematics is inferred from an
unlocated exact title or from combining familiar terminology.
\begin{thebibliography}{99}
\bibitem{nb:clark}
A. Clark, ``The syntactic concept lattice: Another algebraic theory of the
context-free languages?'', \emph{Journal of Logic and Computation},
published online 2013.
\href{https://doi.org/10.1093/logcom/ext037}{doi:10.1093/logcom/ext037}.
\href{https://alexc17.github.io/static/pdfs/JLC2013.pdf}{Author PDF}.

\bibitem{nb:shen-zhang}
L. Shen and D. Zhang, ``The concept lattice functors'',
\emph{International Journal of Approximate Reasoning} 54 (2013), 166--183.
\href{https://doi.org/10.1016/j.ijar.2012.07.002}{doi:10.1016/j.ijar.2012.07.002}.

\bibitem{nb:erne}
M. Ern\'e, ``Categories of contexts'', 2014.
\href{https://arxiv.org/abs/1407.0512}{arXiv:1407.0512}.

\bibitem{nb:gratus-porter}
J. Gratus and T. Porter, ``A geometry of information, I: Nerves, posets
and differential forms'', \emph{Dagstuhl Seminar Proceedings} 04351,
2005.
\href{https://arxiv.org/abs/cs/0512010}{arXiv:cs/0512010}.
\href{https://drops.dagstuhl.de/storage/16dagstuhl-seminar-proceedings/dsp-vol04351/DagSemProc.04351.6/DagSemProc.04351.6.pdf}{Proceedings PDF}.
Relevant statements: Proposition 1 and Corollary 1, pp.~8--9, and the
refinement discussion on pp.~11--12.

\bibitem{nb:brun-salbu}
M. Brun and L. M. Salbu, ``The Rectangle Complex of a Relation'',
\emph{Mediterranean Journal of Mathematics} 20 (2023), article 7.
\href{https://doi.org/10.1007/s00009-022-02213-0}{doi:10.1007/s00009-022-02213-0}.

\bibitem{cg:angluin}
D. Angluin, \emph{Learning Regular Sets from Queries and Counterexamples},
Information and Computation 75 (1987), 87--106.
\href{https://doi.org/10.1016/0890-5401(87)90052-6}{doi:10.1016/0890-5401(87)90052-6}.

\bibitem{cg:universal}
S. Lombardy and J. Sakarovitch, \emph{The Universal Automaton}, in
\emph{Logic and Automata: History and Perspectives}; author version
(2007), Propositions 2.1, 2.6, Theorem 2.16, and Definition 2.4.
\href{https://perso.telecom-paristech.fr/jsaka/PUB/Files/TUA.pdf}{Author PDF}.
Publisher volume records use 2008; proposition numbering is the locator here.

\bibitem{cg:meet}
O. Kl\'{\i}ma and L. Pol\'{a}k, \emph{Syntactic Structures of Regular
Languages}, 2017 version, Section 4.
\href{https://arxiv.org/abs/1612.06247}{arXiv:1612.06247}.

\bibitem{cg:quantaloid}
L. Shen and D. Zhang, \emph{Categories enriched over a quantaloid:
Isbell adjunctions and Kan adjunctions}, Theory and Applications of
Categories 28(20) (2013), 577--615.
\href{https://arxiv.org/abs/1307.5625}{arXiv:1307.5625};
\href{https://emis.de/ft/8978}{Published-paper mirror}.

\bibitem{cg:freund}
A. Freund, M. Andreatta, and J.-L. Giavitto,
\emph{Lattice-based and topological representations of binary relations
with an application to music}, Annals of Mathematics and Artificial
Intelligence 73 (2015), 311--334.
\href{https://doi.org/10.1007/s10472-014-9445-3}{doi:10.1007/s10472-014-9445-3}.

\bibitem{cg:ayzenberg}
A. Ayzenberg, \emph{Topology of nerves and formal concepts} (2019),
Proposition 5.7 and Remark 5.8.
\href{https://arxiv.org/abs/1911.05491}{arXiv:1911.05491}.

\bibitem{cg:dowker-persistence}
S. Chowdhury and F. M\'emoli,
\emph{A functorial Dowker theorem and persistent homology of asymmetric
networks}, Journal of Applied and Computational Topology 2 (2018),
115--175, Theorem 3 and Corollary 20.
\href{https://doi.org/10.1007/s41468-018-0020-6}{doi:10.1007/s41468-018-0020-6}.

\bibitem{cg:selfmap}
H. Edelsbrunner, G. Jab\l o\'nski, and M. Mrozek,
\emph{The Persistent Homology of a Self-Map}, Foundations of Computational
Mathematics 15 (2015), 1213--1244, Section 4.1.
\href{https://doi.org/10.1007/s10208-014-9223-y}{doi:10.1007/s10208-014-9223-y}.

\bibitem{cg:nerve}
U. Bauer, M. Kerber, F. Roll, and A. Rolle,
\emph{A Unified View on the Functorial Nerve Theorem and its Variations},
journal publication 2023; arXiv version 6 (2025), Theorems 3.9, 4.8, 5.9
and Remark 5.12.
\href{https://arxiv.org/abs/2203.03571}{arXiv:2203.03571}.

\bibitem{cg:quillen}
D. Quillen, \emph{Higher Algebraic K-Theory: I}, Lecture Notes in
Mathematics 341 (1973), 85--147, Theorem A.
\href{https://doi.org/10.1007/BFb0067053}{doi:10.1007/BFb0067053}.

\bibitem{cg:green}
R. E. Green, C. A. Joslyn, A. Myers, M. G. Rawson, and M. Robinson,
\emph{The Topological Structures of the Orders of Hypergraphs} (2025),
preprint version 2.
\href{https://arxiv.org/abs/2504.11760}{arXiv:2504.11760}.
\end{thebibliography}
\end{document}
